\section*{Aktivitas Pembelajaran dan Praktikum}
\addcontentsline{toc}{section}{Aktivitas Pembelajaran dan Praktikum}

\begin{table}[H]
\centering
\caption{Rencana aktivitas Bab XVII}
\label{tab:ch17-aktivitas}
\begin{tabularx}{\textwidth}{@{}llX@{}}
\toprule
\textbf{Sesi} & \textbf{Bentuk} & \textbf{Kegiatan} \\
\midrule
1 & Diskusi & Peta pencapaian dan analisis hasil \\
2 & Latihan & Soal UAS dan studi kasus \\
3 & Presentasi & Demo proyek akhir \\
4 & Refleksi & Peer review dan retrospektif \\
\bottomrule
\end{tabularx}
\end{table}

\subsection*{Demo Akhir dan Peer Review}
\textbf{Tujuan.} Mempresentasikan proyek, menerima umpan balik, dan
menyusun rencana perbaikan.

\textbf{Langkah kerja.}
\begin{enumerate}[label=\arabic*.,leftmargin=*]
  \item Siapkan demo 10 menit: arsitektur, fitur, dan keterbatasan.
  \item Tunjukkan satu program valid dan satu program tidak valid.
  \item Jelaskan keputusan desain utama.
  \item Minta teman menguji dengan program baru.
  \item Catat umpan balik dan susun rencana perbaikan.
\end{enumerate}

\textbf{Definisi selesai.}
\begin{checklist}
  \item Demo berjalan tanpa kesalahan.
  \item Program valid dan tidak valid ditunjukkan.
  \item Umpan balik tercatat.
  \item Rencana perbaikan disusun.
\end{checklist}

Demo dan peer review menjadi bagian dari Project Final (20\%).
